% Part: sets-functions-relations-complete

\documentclass[../../include/open-logic-part]{subfiles}

\begin{document}

\olpart{sfr}{Na\"ieve verzamelingenleer}

\begin{editorial}
  Dit deel legt de basis van de naïeve verzamelingenleer uit. Tim Buttons
  Open Set Theory is in dit deel opgenomen. Daardoor legt dit deel ook uit hoe
  getallenstelsels worden opgebouwd en gaat het over oneindigheid. Voor de
  delen van het OLP over logica heb je deze stof niet nodig.
\end{editorial}

\olimport[sets]{sets}

\olimport[relations]{relations-complete}

\olimport[functions]{functions}

\olimport[size-of-sets]{size-of-sets-complete}

\olimport[arithmetization]{arithmetization}

\olimport[infinite]{infinite}

\OLEndPartHook
\end{document}
